% Document CTM1. Canonical source; see codenames.json.
\documentclass[11pt,letterpaper,reqno]{amsart}
% BEGIN INLINE manuscript-style
\makeatletter
\NeedsTeXFormat{LaTeX2e}[2020/10/01]
% Use with \documentclass[11pt,letterpaper,reqno]{amsart}.
% This file controls presentation only. It supplies no mathematical content.

\RequirePackage[T1]{fontenc}
\RequirePackage{lmodern}
\RequirePackage{microtype}
\RequirePackage[letterpaper,margin=1in,headheight=15pt]{geometry}
\RequirePackage{amsmath,amsthm,amssymb,mathtools}
\RequirePackage{enumitem}
\RequirePackage{perpage}
\RequirePackage{xcolor}
\definecolor{manuscriptred}{rgb}{0.6,0.1,0.1}
\definecolor{manuscriptblue}{rgb}{0.1,0.3,0.5}
\definecolor{manuscriptmagenta}{HTML}{60445C}
\RequirePackage[breaklinks=true,colorlinks=true,
  linkcolor=manuscriptred,citecolor=manuscriptblue,urlcolor=manuscriptmagenta]{hyperref}
\RequirePackage{cleveref}
\hypersetup{pdfauthor={}}

% Retain the AMS paragraph, title, heading and running-head conventions.
% Equations retain one global Arabic sequence, as in the reference.
\setcounter{tocdepth}{1}
\renewcommand{\thefigure}{\Roman{figure}}
\renewcommand{\thetable}{\Roman{table}}
\MakePerPage{footnote}
\renewcommand{\thefootnote}{%
  \ifnum\value{footnote}>9\arabic{footnote}\else\fnsymbol{footnote}\fi}

% Upright statements. Main results are A, B, C; supporting statements
% share a section-local sequence. These end marks indicate typography,
% never proof completeness, validation, or a deferred-proof status.
\theoremstyle{definition}
\newtheorem{ms@theorem}{Theorem}
\renewcommand{\thems@theorem}{\Alph{ms@theorem}}
\newtheorem*{ms@theoremstar}{Theorem}
\newtheorem{ms@proposition}{Proposition}[section]
\newtheorem{ms@lemma}[ms@proposition]{Lemma}
\newtheorem{ms@definition}[ms@proposition]{Definition}
\newtheorem{ms@corollary}[ms@proposition]{Corollary}
\newtheorem*{ms@corollarystar}{Corollary}

\newcommand{\ms@majorstart}{%
  \pushQED{\qed}\renewcommand{\qedsymbol}{$\blacksquare\square$}}
\newcommand{\ms@minorstart}{%
  \pushQED{\qed}\renewcommand{\qedsymbol}{$\blacksquare$}}
\newenvironment{theorem}[1][]{\ms@majorstart\ms@theorem[#1]}
  {\popQED\endms@theorem}
\newenvironment{theorem*}[1][]{\ms@majorstart\ms@theoremstar[#1]}
  {\popQED\endms@theoremstar}
\newenvironment{proposition}[1][]{\ms@minorstart\ms@proposition[#1]}
  {\popQED\endms@proposition}
\newenvironment{lemma}[1][]{\ms@minorstart\ms@lemma[#1]}
  {\popQED\endms@lemma}
\newenvironment{definition}[1][]{\ms@minorstart\ms@definition[#1]}
  {\popQED\endms@definition}
\newenvironment{corollary}[1][]{\ms@majorstart\ms@corollary[#1]}
  {\popQED\endms@corollary}
\newenvironment{corollary*}[1][]{\ms@majorstart\ms@corollarystar[#1]}
  {\popQED\endms@corollarystar}
% The source's lemma/proposition variants are deliberately not carried over:
% choose one marker for each statement type rather than encode proof status.

\theoremstyle{remark}
\newtheorem{remark}[ms@proposition]{Remark}
\newtheorem*{remark*}{Remark}
% The ordinary amsthm proof environment retains its open square.

\crefname{ms@theorem}{theorem}{theorems}
\Crefname{ms@theorem}{Theorem}{Theorems}
\crefname{ms@proposition}{proposition}{propositions}
\Crefname{ms@proposition}{Proposition}{Propositions}
\crefname{ms@lemma}{lemma}{lemmas}
\Crefname{ms@lemma}{Lemma}{Lemmas}
\crefname{ms@definition}{definition}{definitions}
\Crefname{ms@definition}{Definition}{Definitions}
\crefname{ms@corollary}{corollary}{corollaries}
\Crefname{ms@corollary}{Corollary}{Corollaries}
\crefname{remark}{remark}{remarks}
\Crefname{remark}{Remark}{Remarks}

% Bibliography configuration for this manuscript.
\makeatother
% END INLINE manuscript-style
\newcommand{\R}{\mathbb R}
\newcommand{\C}{\mathbb C}
\newcommand{\B}{\mathbb B}
\DeclareMathOperator{\supp}{supp}
\newcommand{\dd}{\,\mathrm d}
\title{Invariance at coinciding points}
\date{}
\hypersetup{pdftitle={Invariance at coinciding points},pdfauthor={}}
\begin{document}
\begin{abstract}
A bounded multilinear form on vector-valued \(L^2(\R^3)\) vanishes
if it vanishes on entries with disjoint planar supports and whenever
one entry is a compactly supported gradient. We prove this statement
and apply it to moments of closed, Gaussian-dominated two-form fields.
It extends rotation identities from separated planar supports to all
\(L^2\) tests and hence to the law of the field.
\end{abstract}
\maketitle
\tableofcontents

\section{Introduction}\label{sec:introduction}
Write $x=(y,z)\in\R^2\times\R$, and, for an open set $U\subset\R^2$, put
\[
 \mathcal L_U=\{f\in L^2(\R^3;\C^3):
       f=0\text{ almost everywhere outside }U\times\R\}.
\]
All norms on these spaces are $L^2$ norms. We say that a field lies over
$K\subset\R^2$ when it vanishes almost everywhere outside $K\times\R$.
Separation therefore concerns the projected planar supports; the
transverse coordinate $z$ remains smeared.

\begin{theorem}[Absence of contact terms]\label{thm:contact}
Let $U\subset\R^2$ be open and $N\ge1$. Suppose that
$D\colon\mathcal L_U^N\to\C$ is a complex multilinear form satisfying
\[
 |D(f_1,\ldots,f_N)|\le M\prod_{a=1}^N\|f_a\|_2
       \qquad(f_a\in\mathcal L_U)
\]
for some $M<\infty$. Assume that:
\begin{enumerate}[label=\textup{(\roman*)}]
\item $D(f_1,\ldots,f_N)=0$ if the entries lie over compact sets
$K_1,\ldots,K_N\subset U$ that are pairwise disjoint;
\item $D(f_1,\ldots,f_N)=0$ whenever one entry is $\nabla u$ for some
$u\in C_c^\infty(U\times\R;\C)$.
\end{enumerate}
Then $D=0$.
\end{theorem}

An identity for correlation functions at separated arguments need not
determine the same identity at coinciding arguments. In a
Gaussian-dominated field theory the moments are bounded multilinear
forms on \(L^2\), so their values are defined even when the test supports
overlap. The gradient hypothesis above is the additional condition
supplied by closedness of a two-form in three dimensions.

The theorem is Proposition~8.4 of \cite{QEDCurrent}.
No probability law enters its hypotheses.
Its application to rotations is direct: the difference between a rotated and
an unrotated moment is bounded, and closedness makes it vanish on gradients.
The remaining hypothesis is precisely the invariance already obtained away
from coinciding planar arguments.

The proof has two different steps. For a bilinear form, locality forces its
representing operator to commute with multiplication by functions of $y$.
One can then multiply gradients by arbitrary bounded planar functions and
obtain a total family of vector fields. For three or more entries, induction
removes terms involving disjoint groups of supports. A partition into small
squares leaves only terms for which all entries lie over the same square;
their sum tends to zero by the $L^2$ bound. Both steps are needed. Boundedness
alone does not remove bilinear forms supported on the diagonal.

\section{Locality and gradients}

We prove \cref{thm:contact} simultaneously for all open sets $U$, by
induction on $N$. A bounded Borel function $h$ on $\R^2$ acts on
$\mathcal L_U$ by
\[
 (hf)(y,z)=h(y)f(y,z).
\]
For a Borel set $A$, the same notation $1_A$ denotes its indicator and the
corresponding orthogonal projection. Fields lying over compact subsets of
$U$ are dense in $\mathcal L_U$: take an increasing compact exhaustion of
$U$ and use dominated convergence. When $N=1$, the first hypothesis of the
theorem therefore gives the conclusion.

For $N=2$, apply the Hilbert-space representation theorem
\cite[Theorem 3.3]{Melrose} to the bounded
sesquilinear form $(f',f)\mapsto D(\overline{f'},f)$. There is a bounded
operator $T$ on $\mathcal L_U$ such that
\begin{equation}\label{eq:represent}
 D(f_1,f_2)=\int_{U\times\R}f_1(x)\cdot(Tf_2)(x)\dd x,
\end{equation}
where the dot denotes the complex bilinear pairing on $\C^3$.

The first hypothesis implies $1_KT1_{K'}=0$ for disjoint compact subsets
$K,K'$ of $U$. Indeed, its matrix elements in \eqref{eq:represent} are
$D(1_Kf_1,1_{K'}f_2)$, which vanish. The same conclusion holds for disjoint
Borel subsets $A,A'$ of $U$. To see this, use inner regularity of planar
Lebesgue measure on the bounded sets $A\cap\{|y|<n\}$ and
$A'\cap\{|y|<n\}$ to choose increasing compact sets $K_n\subset A$ and
$K_n'\subset A'$ whose unions exhaust $A$ and $A'$ up to null sets. Then
$1_{K_n}\to1_A$ and $1_{K_n'}\to1_{A'}$ strongly on $\mathcal L_U$, so
\[
 1_AT1_{A'}=\lim_{n\to\infty}1_{K_n}T1_{K_n'}=0
\]
in the strong operator topology.

Take $A'=U\setminus A$ and then interchange the two sets. The resulting
identities give
\[
 T1_A=1_AT1_A=1_AT.
\]
Since every bounded Borel function is a uniform limit of simple Borel
functions, $T$ commutes with all planar multipliers. Consequently
\begin{equation}\label{eq:balance}
 D(hf_1,f_2)=D(f_1,hf_2).
\end{equation}
The gradient hypothesis and this identity imply
\begin{equation}\label{eq:modulated}
 D(h\nabla u,f)=0,
 \qquad u\in C_c^\infty(U\times\R;\C),\quad f\in\mathcal L_U,
\end{equation}
for every bounded Borel function $h$ on $\R^2$.

We now use \eqref{eq:modulated} to separate the three components of a vector
field. Fix a compact set $A\subset U$, choose $u_0\in C_c^\infty(U)$ equal
to one near $A$, and let $\chi\in C_c^\infty(\R)$. Write $y=(y^0,y^1)$.
The gradients of
\[
 u_0(y)\chi(z),\qquad
 u_0(y)y^0\chi(z),\qquad
 u_0(y)y^1\chi(z)
\]
restrict on $A\times\R$ to
\[
 \chi'(z)e_3,\qquad
 \chi(z)e_1+y^0\chi'(z)e_3,\qquad
 \chi(z)e_2+y^1\chi'(z)e_3,
\]
respectively. Apply \eqref{eq:modulated} with the multiplier $h1_A$.
The first gradient gives
\[
 D(h1_A\chi'e_3,f)=0.
\]
In the second identity the term containing $e_3$ vanishes by this formula,
with the bounded multiplier $hy^01_A$. The third identity is treated with
$hy^11_A$. Thus
\begin{equation}\label{eq:three-families}
 D(h1_A\chi e_1,f)=D(h1_A\chi e_2,f)
                  =D(h1_A\chi'e_3,f)=0.
\end{equation}

The fields in \eqref{eq:three-families}, as $A,h,\chi$ vary, have dense
linear span in $\mathcal L_U$. For the first two components this is the
density of products of planar functions of compact support and smooth
compactly supported functions of $z$. For the third component, the
derivatives $\chi'$ are themselves total in $L^2(\R)$. A function
orthogonal to every such derivative has zero distributional derivative,
and hence is almost everywhere constant
\cite[Proposition 1.4.8]{Gerard}; the only constant in $L^2(\R)$
is zero. It follows that the continuous form $D(\cdot,f)$ vanishes on all
of $\mathcal L_U$. This proves the theorem for two entries.

\section{Several coinciding arguments}

Let $N\ge3$, and suppose that the theorem holds for every smaller number
of entries and every open subset of $\R^2$. Before partitioning the
supports, we record two consequences of the induction hypothesis.

First, fix an index $i$ and a compact set $K\subset U$. If $f_i$ lies over
$K$ and all the other entries belong to $\mathcal L_{U\setminus K}$,
then
\begin{equation}\label{eq:one-separated}
 D(f_1,\ldots,f_N)=0.
\end{equation}
Indeed, hold $f_i$ fixed. The remaining entries define a bounded
$(N-1)$-linear form on $\mathcal L_{U\setminus K}$. It vanishes on
gradients there, because a compactly supported smooth function on
$(U\setminus K)\times\R$, extended by zero, belongs to
$C_c^\infty(U\times\R;\C)$.
It also vanishes on entries over pairwise disjoint compact subsets of
$U\setminus K$: those compacts, together with $K$, satisfy the first
hypothesis for $D$. Induction gives \eqref{eq:one-separated}.

Second, partition the indices into two nonempty sets $I,J$, and let
$K_J\subset U$ be compact. If $f_j$ lies over $K_J$ for every $j\in J$
and $f_i\in\mathcal L_{U\setminus K_J}$ for every $i\in I$, then
\begin{equation}\label{eq:clusters-separated}
 D(f_1,\ldots,f_N)=0.
\end{equation}
Hold the entries in $J$ fixed and regard $D$ as a form with $|I|<N$
entries on $\mathcal L_{U\setminus K_J}$. The gradient hypothesis is
inherited. To check the other hypothesis, suppose that the entries with
indices in $I$ lie over pairwise disjoint compact sets $K_i$ in
$U\setminus K_J$. Choose one $i'\in I$. Every other entry lies outside
$K_{i'}\times\R$, so \eqref{eq:one-separated}, with $K=K_{i'}$, gives
vanishing. Induction applied to this $|I|$-linear form proves
\eqref{eq:clusters-separated}.

By density and continuity, it remains to treat entries that all lie over
one compact set $K_0\subset U$. For $d>0$, let
\[
 \mathcal Q_d=\{d(k+[0,1)^2):k\in\mathbb Z^2\}.
\]
Only finitely many of these half-open squares meet $K_0$. Multilinearity
therefore gives
\begin{equation}\label{eq:partition}
 D(f_1,\ldots,f_N)=
 \sum_{Q_1,\ldots,Q_N\in\mathcal Q_d}
 D(1_{Q_1}f_1,\ldots,1_{Q_N}f_N).
\end{equation}
The sum has only finitely many nonzero terms. We claim that every term
with squares $Q_a$ not all equal vanishes.

For such a tuple put $I=\{i:Q_i=Q_1\}$ and let $J$ be its nonempty
complement. For $0<\varepsilon<d/2$, denote by $Q^\varepsilon$ the
closed square of side $d-2\varepsilon$ concentric with $Q$.
The compact sets
\[
 K_0\cap Q_1^\varepsilon,
 \qquad K_J=K_0\cap\bigcup_{j\in J}Q_j^\varepsilon
\]
are disjoint. The entries $1_{Q_j^\varepsilon}f_j$, $j\in J$, lie
over $K_J$, and the entries with indices in $I$ lie outside $K_J$.
Equation~\eqref{eq:clusters-separated} gives
\[
 D(1_{Q_1^\varepsilon}f_1,\ldots,
                         1_{Q_N^\varepsilon}f_N)=0.
\]
For each $a$, $1_{Q_a^\varepsilon}f_a\to1_{Q_a}f_a$ in $L^2$ as
$\varepsilon\downarrow0$, since the boundary of a square has measure
zero. Continuity proves the claim.

Only the terms in which all entries lie over the same square remain in
\eqref{eq:partition}. Their sum satisfies
\begin{align}
 |D(f_1,\ldots,f_N)|
 &\le M\sum_{Q\in\mathcal Q_d}\prod_{a=1}^N\|1_Qf_a\|_2\notag\\
 &\le M\left(\max_{Q\in\mathcal Q_d}\|1_Qf_3\|_2\right)
             \prod_{a\ne3}\|f_a\|_2.
 \label{eq:small-square}
\end{align}
For the last inequality use Cauchy--Schwarz in the sum for the first
two entries and $\|1_Qf_a\|_2\le\|f_a\|_2$ for $a\ge4$.
The function
\[
 y\longmapsto\int_\R|f_3(y,z)|^2\dd z
\]
belongs to $L^1(\R^2)$ by Tonelli's theorem. Each square has area $d^2$,
so absolute continuity of its integral gives
\[
 \max_{Q\in\mathcal Q_d}\|1_Qf_3\|_2\longrightarrow0
                  \qquad(d\downarrow0).
\]
The left side of \eqref{eq:small-square} is independent of $d$.
It is therefore zero. This completes the induction and the proof of
\cref{thm:contact}.\qed

\section{The rotation defect}

For $g\in O(3)$ we use the action on pseudovector fields
\begin{equation}\label{eq:action}
 (g\cdot b)(x)=\det(g)\,g\,b(g^{-1}x).
\end{equation}
This is the action obtained from two-forms under the Hodge isometry
$*f=(f_{23},-f_{13},f_{12})$. Let $R_t$ be rotation by $t$ in the
$(x_1,x_2)$-plane, fixing $x_3$.

\begin{corollary}[From separated supports to all tests]\label{cor:rotation}
Let $N\ge1$, and let $S_N$ be a bounded complex $N$-linear form on
$L^2(\R^3;\C^3)$. Assume that it vanishes whenever one entry is
$\nabla u$, $u\in C_c^\infty(\R^3;\C)$. Fix $t\in\R$, and suppose that
\[
 S_N(R_t\cdot b_1,\ldots,R_t\cdot b_N)=S_N(b_1,\ldots,b_N)
\]
whenever the entries lie over pairwise disjoint compact subsets of
$\R^2$. Then the same identity holds for all entries in
$L^2(\R^3;\C^3)$.
\end{corollary}

\begin{proof}
Apply \cref{thm:contact}, with $U=\R^2$, to
\[
 D_t(b_1,\ldots,b_N)
 =S_N(R_t\cdot b_1,\ldots,R_t\cdot b_N)-S_N(b_1,\ldots,b_N).
\]
The action of $R_t$ is an $L^2$ isometry, so $D_t$ is bounded.
The first vanishing hypothesis is the assumed separated-support
identity. For the second, use
\[
 R_t\cdot\nabla u=\nabla(u\circ R_{-t}).
\]
Since $u\circ R_{-t}\in C_c^\infty(\R^3;\C)$, both terms defining
$D_t$ vanish when an entry is $\nabla u$.
\end{proof}

For the application to curvature, let $\mu$ be a probability law of
real tempered two-forms $F$ on $\R^3$, with the sigma-algebra generated
by evaluations on real Schwartz tests, and put
$\B(b)=F(*^{-1}b)$ on Schwartz vector fields. Gaussian domination means
\begin{equation}\label{eq:G}
 \int e^{\B(b)}\dd\mu\le e^{\|b\|_2^2/2}
           \qquad(b\in\mathcal S(\R^3;\R^3)).
\end{equation}
Closedness means $dF=0$ almost surely. In vector notation it gives
\begin{equation}\label{eq:gradient}
 \B(\nabla u)=0
       \qquad(u\in C_c^\infty(\R^3;\C))
\end{equation}
almost surely; see \cite[Section~8.5]{QEDCurrent} for this
vector form of closedness.

By \cite[Proposition~5.1]{QEDCurrent}, the map $b\mapsto\B(b)$
extends to a continuous complex linear map from $L^2$ to every
$L^p(\mu)$, $1\le p<\infty$. Its moments
\[
 S_N(b_1,\ldots,b_N)=\int\B(b_1)\cdots\B(b_N)\dd\mu
\]
are bounded symmetric multilinear forms. More precisely,
\[
 \|\B(b)\|_{L^N(\mu)}\le c_N\|b\|_2,
 \qquad
 |S_N(b_1,\ldots,b_N)|\le c_N^N\prod_a\|b_a\|_2,
\]
with constants $c_N$ depending only on $N$.
The moment bound follows from the field bound by H\"older's inequality.
In particular no condition on the relative supports is needed to define
$S_N$. Equation~\eqref{eq:gradient} gives the gradient hypothesis of
\cref{cor:rotation}.

\begin{proposition}[Invariance of a closed Gaussian-dominated law]
\label{prop:law}
Suppose that $\mu$ satisfies \eqref{eq:G} and is almost surely closed.
Assume that, for every $N\ge1$ and every $t\in\R$, its moments satisfy
the rotation identity of \cref{cor:rotation} on entries over
pairwise disjoint compact planar sets. Then all moments satisfy that
identity on $L^2(\R^3;\C^3)$, and $(R_t)_*\mu=\mu$ for every $t$.
\end{proposition}

\begin{proof}
The moment bounds and \eqref{eq:gradient} verify the hypotheses of
\cref{cor:rotation}, which proves the first assertion.
To recover the law, take a real Schwartz field $b$. The partial sums
of the exponential series for $e^{i\B(b)}$ have modulus at most
\[
 e^{|\B(b)|}\le e^{\B(b)}+e^{-\B(b)}.
\]
The last expression is integrable by \eqref{eq:G}. Dominated convergence
therefore gives
\[
 \int e^{i\B(b)}\dd\mu
 =1+\sum_{N\ge1}\frac{i^N}{N!}S_N(b,\ldots,b).
\]
The series for $R_t\cdot b$ has the same terms by the first assertion.
The characteristic functional is thus invariant. Laws on tempered
distributions are determined by their characteristic functionals,
so $(R_t)_*\mu=\mu$; see also
\cite[Section~5.2]{QEDCurrent}.
\end{proof}

This is the final step in the proof of
\cite[Theorem~8.6]{QEDCurrent}. For that application, angular
covariance of the plate kernels and their identification as moment
densities on all distinct configurations
\cite[Sections~8.9--8.10]{QEDCurrent} give the separated-support
identity. \Cref{thm:contact} then uses the bounded moments and their
gradient identities to remove the support restriction.

\section{What the domain specifies at coincidence}

The term ``invariance at coinciding arguments'' must be read through
the domain of the correlation functional. For a distribution defined
on the open configuration space
\[
 (\R^d)^N_{\ne}
 =\{(x_1,\ldots,x_N):x_a\ne x_b\text{ for }a\ne b\},
\]
covariance is an identity on that open set. It gives no values of an
extension on the diagonals. A formulation on Schwartz functions flat
at every diagonal likewise imposes its identities on that stated
test-function space. A functional given only on separated compactly
supported $C^1$ tests likewise specifies no values on overlapping
supports. This domain should be kept distinct from the unrestricted
$L^2$ moments used here.

In the Gaussian-dominated three-dimensional setting, overlapping
supports are already allowed before invariance is proved. It remains
to extend the covariance identity from tests with disjoint planar
supports to all $L^2$ tests: the analytic plate kernels are first
identified at pairwise distinct planar points.
The corresponding rotation defect may still be supported where some
planar arguments coincide. \Cref{thm:contact} excludes that
defect by combining the $L^2$ bound with the gradient identities.
It does not assign a pointwise value to a singular correlation
kernel on the diagonal. Its conclusion is an identity between
bounded multilinear forms for all $L^2$ test fields.

\begin{thebibliography}{9}
\bibitem{Gerard} P.~G\'erard, \emph{Distributions and Partial Differential
Equations}, lecture notes, Fall 2019, December 12, 2019.
\href{https://www.imo.universite-paris-saclay.fr/~patrick.gerard/Distributions2019_Chap1,2,3.pdf}{Lecture notes}.
\bibitem{Melrose} R.~Melrose, \emph{Functional Analysis: Lecture notes
for 18.102, Spring 2020}, Massachusetts Institute of Technology, 2020.
\href{https://ocw.mit.edu/courses/18-102-introduction-to-functional-analysis-spring-2021/3d4cc88026d44a01f936cd6a0aa995cb_MIT18_102s20_lec_FA.pdf}{Lecture notes}.
\bibitem{QEDCurrent}
\emph{The field-strength Osterwalder--Schrader axioms for QED
in three dimensions}, comprehensive edition, 263~pp., 2026.
\url{https://github.com/ethanwsussman/Yang--Mills-LLM-experiment/blob/main/QED3/qed3.pdf}.
\end{thebibliography}
\end{document}
